\section{Phân công công việc}

\subsection{Kế hoạch phân công}
\label{subsec:ke-hoach-phan-cong}

Kế hoạch dưới đây cụ thể hóa phạm vi phụ trách của từng thành viên trong
Version 1. Nội dung tập trung vào phân tích yêu cầu, thiết kế giao diện,
thiết kế cơ sở dữ liệu và chuẩn bị hồ sơ nộp. Đây là phân công dự kiến,
không phải xác nhận các công việc đã hoàn thành.

Mã \texttt{F01}, \texttt{F02}, \ldots\ nhận biết chức năng; mã \texttt{T01},
\texttt{T02}, \ldots\ nhận biết công việc. Phân hệ được thể hiện trong danh mục
chức năng tại mục~\ref{subsubsec:client-user} và
mục~\ref{subsubsec:client-admin}, không đưa vào mã.
Mỗi mã được giữ ổn định khi đổi tên hoặc sắp xếp lại bảng.

Một công việc có thể liên quan đến nhiều chức năng và một chức năng có thể
được đối chiếu với nhiều công việc. Trong V1, liên kết này thể hiện phạm vi
thiết kế và kiểm tra, không có nghĩa chức năng đã được triển khai.

\begingroup
\small
\setlength{\tabcolsep}{4pt}
\begin{longtable}{|C{1cm}|L{\dimexpr\textwidth-6.6cm-8\tabcolsep-5\arrayrulewidth\relax}|L{3.2cm}|C{2.4cm}|}
  \caption{Kế hoạch phân công công việc trong Version 1}
  \label{tab:ke-hoach-phan-cong}\\
  \hline
  \textbf{Mã} & \textbf{Công việc} & \textbf{Người phụ trách} &
  \textbf{Chức năng liên quan}\\
  \hline
  \endfirsthead
  \hline
  \textbf{Mã} & \textbf{Công việc} & \textbf{Người phụ trách} &
  \textbf{Chức năng liên quan}\\
  \hline
  \endhead
  \hypertarget{T01}{\texttt{T01}} &
    Phân tích đề bài, xác định phạm vi và các yêu cầu của Version 1.
    & Cả hai thành viên & \hyperlink{F01}{F01}--\hyperlink{F21}{F21}\\
  \hline
  \hypertarget{T02}{\texttt{T02}} &
    Thiết kế giao diện phân hệ quản trị viên, bao gồm các biểu mẫu,
      danh sách, điều khiển lọc, sắp xếp, thống kê và biểu đồ theo đề bài.
    & Châu Vũ Trung & \hyperlink{F13}{F13}--\hyperlink{F21}{F21}\\
  \hline
  \hypertarget{T03}{\texttt{T03}} &
    Thiết kế giao diện phân hệ người sử dụng, bao gồm tài khoản,
      bạn bè, tìm người, chat riêng, chat nhóm và các thao tác liên quan.
    & Nguyễn Sơn Hải & \hyperlink{F01}{F01}--\hyperlink{F12}{F12}\\
  \hline
  \hypertarget{T04}{\texttt{T04}} &
    Thiết kế cơ sở dữ liệu và chuẩn bị script tạo cơ sở dữ liệu PostgreSQL.
    & Cả hai thành viên & \hyperlink{F01}{F01}--\hyperlink{F21}{F21}\\
  \hline
  \hypertarget{T05}{\texttt{T05}} &
    Đối chiếu các màn hình và thiết kế cơ sở dữ liệu với yêu cầu đề bài,
      ghi nhận các nội dung cần bổ sung hoặc điều chỉnh.
    & Cả hai thành viên & \hyperlink{F01}{F01}--\hyperlink{F21}{F21}\\
  \hline
  \hypertarget{T06}{\texttt{T06}} &
    Chuẩn bị nội dung báo cáo và ảnh minh họa của phân hệ quản trị viên.
    & Châu Vũ Trung & \hyperlink{F13}{F13}--\hyperlink{F21}{F21}\\
  \hline
  \hypertarget{T07}{\texttt{T07}} &
    Chuẩn bị nội dung báo cáo và ảnh minh họa của phân hệ người sử dụng.
    & Nguyễn Sơn Hải & \hyperlink{F01}{F01}--\hyperlink{F12}{F12}\\
  \hline
  \hypertarget{T08}{\texttt{T08}} &
    Tổng hợp báo cáo, thống nhất cách trình bày và sắp xếp hồ sơ nộp V1.
    & Châu Vũ Trung & Công việc chung\\
  \hline
  \hypertarget{T09}{\texttt{T09}} &
    Commit phần công việc thực hiện và cung cấp ảnh commit cá nhân
      để đưa vào hồ sơ nộp.
    & Mỗi thành viên & Công việc chung\\
  \hline
\end{longtable}
\endgroup

Khoảng mã như \hyperlink{F01}{F01}--\hyperlink{F12}{F12} bao gồm tất cả
chức năng từ F01 đến F12; mỗi đầu mút liên kết đến dòng chức năng tương ứng.
Các công việc chung về tổng hợp và quản lý sản phẩm không gắn ép với
một chức năng cụ thể.

\noindent\textbf{Nguyên tắc nghiệm thu:}
Mọi công việc đều phải được cả hai thành viên kiểm tra và thống nhất
công nhận trước khi ghi nhận hoàn thành.

Phân công có thể được điều chỉnh theo khối lượng thực tế sau khi hai thành viên
thống nhất. Kết quả thực hiện được cập nhật tại mục~\ref{subsec:ket-qua-thuc-hien}
và là căn cứ để xác định đóng góp của từng thành viên.

\clearpage
\subsection{Kết quả thực hiện}
\label{subsec:ket-qua-thuc-hien}

Bảng dưới đây ghi nhận công việc thực tế có sản phẩm và minh chứng đối chiếu.
Ngày kiểm tra kỹ thuật và ngày ghi nhận từ commit được chú thích trong bảng.
Ngày hoàn thành theo nghiệm thu chung do hai thành viên xác nhận.

Cột mã công việc sử dụng lại mã \hyperlink{T01}{T01}--\hyperlink{T09}{T09}
ở mục~\ref{subsec:ke-hoach-phan-cong}, không đánh số mới theo
thứ tự ghi nhận. Một công việc có thể có nhiều dòng kết quả nếu được thực hiện
qua nhiều đợt; các chức năng liên quan được đối chiếu từ mã công việc đó.

Ngày hoàn thành được ghi theo định dạng \textit{dd/mm/yyyy}, tương ứng với thời điểm
công việc được cả hai thành viên kiểm tra và công nhận. Nếu chưa có ngày xác nhận,
ngày ghi nhận từ commit được chú thích trong bảng. Minh chứng là đường dẫn
đến sản phẩm hoặc mã commit phù hợp: mã nguồn và script SQL dẫn đến
\href{https://github.com/nguyenshai/Project-java}{GitHub của nhóm};
tài liệu, ảnh và báo cáo dẫn đến
\href{https://drive.google.com/drive/folders/1Kci4Ug_HXLRbBBFc0h6lVz0RbuvbcJTF?usp=sharing}
{Google Drive của nhóm}.

\begin{center}
\small
\setlength{\tabcolsep}{4pt}
\captionof{table}{Kết quả thực hiện công việc}
\label{tab:ket-qua-thuc-hien}
\begin{tabularx}{\textwidth}{|C{1.4cm}|L{2.2cm}|Y|L{2.1cm}|L{2.5cm}|L{2cm}|}
  \hline
  \textbf{Mã công việc} & \textbf{Người thực hiện} & \textbf{Công việc đã thực hiện} &
  \textbf{Ngày hoàn thành} & \textbf{Kết quả / sản phẩm} & \textbf{Link minh chứng}\\
  \hline
  \hyperlink{T08}{\texttt{T08}} & Châu Vũ Trung &
    Thiết lập cấu trúc và bản nháp báo cáo V1 bằng LaTeX. &
    05/11/2026 (commit) & Mã nguồn báo cáo V1 &
    \href{https://github.com/nguyenshai/Project-java/commit/97c5565ffd339990c2cdc7d8f282d2d1a080a7d6}{\texttt{97c5565}} \\
  \hline
  \hyperlink{T08}{\texttt{T08}} & Châu Vũ Trung &
    Thiết lập Maven Wrapper và cấu trúc dự án Java gồm bốn module. &
    05/11/2026 (commit) & Khung dự án Maven &
    \href{https://github.com/nguyenshai/Project-java/commit/97c5565ffd339990c2cdc7d8f282d2d1a080a7d6}{\texttt{97c5565}} \\
  \hline
  \hyperlink{T02}{\texttt{T02}} & Châu Vũ Trung &
    Thiết kế Figma và dựng giao diện quản trị bằng Java Swing cho F13--F21. &
    11/11/2026 (kiểm tra) & Nguyên mẫu và mã nguồn Admin &
    \href{https://www.figma.com/design/maCogjZs9eOOhfkxsOZoIU/ChatApp-Admin?node-id=2-80}{Figma};
    \href{https://github.com/nguyenshai/Project-java/tree/trung/src/admin-client}{mã nguồn} \\
  \hline
  \hyperlink{T06}{\texttt{T06}} & Châu Vũ Trung &
    Viết mục Admin, đối chiếu F13--F21 và kết xuất ảnh từ ứng dụng Swing. &
    11/11/2026 (kiểm tra) & Nội dung và ảnh giao diện Admin &
    Mục~\ref{subsubsec:client-admin}, Hình~\ref{fig:admin-swing-users-history} \\
  \hline
  \rule{0pt}{1.15cm} & & & & & \\
  \hline
  \rule{0pt}{1.15cm} & & & & & \\
  \hline
\end{tabularx}
\end{center}

\clearpage
\subsection{Bảng đóng góp của các thành viên}
\label{subsec:dong-gop-thanh-vien}

Tỷ lệ đóng góp được xác định dựa trên khối lượng công việc thực tế, mức độ
hoàn thành, chất lượng sản phẩm và sự tham gia vào các công việc chung.
Nội dung và tỷ lệ đóng góp được hai thành viên thống nhất trước khi nộp
báo cáo, với tổng tỷ lệ bằng 100\%.

\begin{center}
\small
\setlength{\tabcolsep}{4pt}
\captionof{table}{Đóng góp của các thành viên}
\label{tab:dong-gop-thanh-vien}
\begin{tabularx}{\textwidth}{|C{0.7cm}|L{3cm}|C{2.1cm}|Y|C{2cm}|}
  \hline
  \textbf{STT} & \textbf{Họ và tên} & \textbf{MSSV} &
  \textbf{Nội dung đóng góp thực tế} & \textbf{Tỷ lệ đóng góp}\\
  \hline
  \rule{0pt}{1.1cm}1 & Nguyễn Sơn Hải & 24127163 & & \\
  \hline
  \rule{0pt}{1.1cm}2 & Châu Vũ Trung & 24127256 &
  Thiết lập khung Maven và bản nháp báo cáo; thiết kế, dựng và kiểm tra giao diện Admin F13--F21;
  viết mục Admin và chuẩn bị ảnh giao diện. & \\
  \hline
  \multicolumn{4}{|r|}{\textbf{Tổng cộng}} & \textbf{100\%}\\
  \hline
\end{tabularx}
\end{center}

Nội dung đóng góp và tỷ lệ của từng thành viên sẽ được cập nhật sau khi hoàn thành
các công việc trong V1. Minh chứng được đối chiếu với kết quả tại
mục~\ref{subsec:ket-qua-thuc-hien}, lịch sử commit trên GitHub và tài liệu
trên Google Drive; số lượng commit không phải tiêu chí đánh giá duy nhất.
